\DocumentMetadata{
  pdfversion=2.0,
  pdfstandard=ua-2,
  lang=en-US,
  tagging=on,
  tagging-setup={math/setup=mathml-SE,tabsorder=structure}
}
\documentclass[11pt,letterpaper]{article}
\usepackage[margin=0.68in,includefoot]{geometry}
\usepackage{newcomputermodern}
\usepackage{amsmath,amscd,mathtools}
\usepackage{tikz,adjustbox}
\providecommand{\square}{\mdlgwhtsquare}
\usepackage[hidelinks]{hyperref}
\hypersetup{
  pdftitle={Analysis PhD Exam, January 2017},
  pdfauthor={University of Florida Department of Mathematics},
  pdfsubject={Graduate mathematics examination},
  pdfdisplaydoctitle=true,
  bookmarksopen=true
}
\setcounter{secnumdepth}{0}
\setlength{\parindent}{0pt}
\setlength{\parskip}{0.28em}
\setlength{\emergencystretch}{3em}
\setlength{\leftmargini}{2.15em}
\setlength{\leftmarginii}{2.65em}
\setlength{\leftmarginiii}{3em}
\pagestyle{empty}
\raggedbottom
\allowdisplaybreaks
\NewTaggingSocketPlug{block/list/label}{examproblem}{%
  \tagstructbegin{tag=Lbl}\tagstructend%
  \tagstructbegin{tag=\LBody}%
  \tagstructbegin{tag=\ProblemHeadingTag,title-o={\ProblemHeadingText}}%
  \tagmcbegin{tag=\ProblemHeadingTag,actualtext={\ProblemHeadingText}}%
  #2%
  \tagmcend%
  \tagstructend%
}
\newenvironment{examproblems}[2]{%
  \def\ProblemHeadingTag{#2}%
  \AssignTaggingSocketPlug{block/list/label}{examproblem}%
  \begin{enumerate}%
  \setcounter{enumi}{#1}%
  \renewcommand{\labelenumi}{\textbf{\theenumi.}}%
  \setlength{\topsep}{0.35em}%
  \setlength{\partopsep}{0pt}%
  \setlength{\itemsep}{0.48em}%
  \setlength{\parsep}{0.18em}%
}{\end{enumerate}}
\newenvironment{examsubparts}{%
  \AssignTaggingSocketPlug{block/list/label}{default}%
  \begin{enumerate}%
  \setlength{\topsep}{0.18em}%
  \setlength{\partopsep}{0pt}%
  \setlength{\itemsep}{0.16em}%
  \setlength{\parsep}{0.1em}%
}{\end{enumerate}}
\newenvironment{examinstructions}{%
  \par\begingroup\itshape
}{%
  \par\endgroup\vspace{0.18em}
}
\makeatletter
\renewcommand\subsection{\@startsection{subsection}{2}{0pt}%
  {-1.25ex plus -.25ex minus -.1ex}%
  {0.55ex plus .1ex}%
  {\normalfont\large\bfseries}}
\renewcommand{\@maketitle}{%
  \newpage\null\vskip -1em%
  {\footnotesize\bfseries
    \href{https://www.ufl.edu/}{University of Florida}, \href{https://math.ufl.edu/}{Department of Mathematics}\par}%
  \vskip -0.5em%
  \tagstructbegin{tag=H1,title={Analysis PhD Exam, January 2017}}%
  \tagpdfparaOff%
  \tagmcbegin{tag=H1}%
    {\LARGE\bfseries\href{https://gma.math.ufl.edu/exams/analysis-phd/}{Analysis PhD Exam}, January 2017\par}%
  \tagmcend%
  \tagpdfparaOn%
  \tagstructend%
  \vskip 0.25em%
  \tagmcbegin{artifact}%
    \hrule height 1pt%
  \tagmcend%
  \vskip 1em%
}
\makeatother
\title{Analysis PhD Exam, January 2017}
\author{}
\date{}
\begin{document}
\maketitle
\thispagestyle{empty}
\begin{examinstructions}
Do six of eight. Write solutions in a neat and logical fashion, giving complete reasons for all steps.
\end{examinstructions}
\begin{examproblems}{0}{H2}
\def\ProblemHeadingText{Problem 1.}
\item
Suppose \((X,\mathcal M)\) and \((Y,\mathcal N)\) are measurable spaces, \(\mathcal E\subset\mathcal N\) and \(f:X\to Y\). Show, if the~\(\sigma\)-algebra generated by \(\mathcal E\) contains \(\mathcal N\) and \(f^{-1}(E)\in\mathcal M\) for all \(E\in\mathcal E\), then~\(f\) is \(\mathcal M\)–\(\mathcal N\) measurable.
\def\ProblemHeadingText{Problem 2.}
\item
Suppose \((X,\mathcal M,\mu)\) is a measure space, \(f:X\to[0,\infty)\) is measurable and \(\int_X f<\infty\). For positive integers~\(n\), let \(E_n=\{f\leq n\}\). Show,
\begin{examsubparts}
\item[\textnormal{(a)}]
the sequence \(\left(\int_{E_n}f\right)\) converges to \(\int f\);
\item[\textnormal{(b)}]
for each \(\epsilon>0\) there is a positive integer~\(M\) such that
\[
\int_{E_M^c}f<\epsilon;
\]
\item[\textnormal{(c)}]
for each \(\epsilon>0\) there is a \(\delta>0\) such that if \(E\in\mathcal M\) and \(\mu(E)<\delta\), then \(\int_E f<\epsilon\).
\end{examsubparts}
\def\ProblemHeadingText{Problem 3.}
\item
Let \((X,\mathcal M,\mu)\) be a measure space with completion \((X,\overline{\mathcal M},\overline\mu)\). Show, if \(f:X\to\mathbb R\) is \(\overline\mu\) measurable, then there exists a~\(\mu\) measurable \(g:X\to\mathbb R\) such that \(\overline\mu(\{x\in X:f(x)\neq g(x)\})=0\).
\def\ProblemHeadingText{Problem 4.}
\item
Given \(f\in L^1([0,1])\), let \(g(x)=\int_x^1\frac{f(t)}{t}\,dt\). Show \(g\in L^1([0,1])\) and
\[
\int_0^1g(x)\,dx=\int_0^1f(t)\,dt.
\]
\def\ProblemHeadingText{Problem 5.}
\item
Let~\(X\) be a normed vector space (over \(\mathbb C\)). Suppose \(\mathcal M\subset X\) is a closed subspace and \(x\in X\setminus\mathcal M\). Prove the subspace \(\mathcal N=\mathcal M+\mathbb Cx\) is closed in~\(X\). Prove if \(\mathcal L\subset X\) is a finite dimensional subspace, then \(\mathcal L\) is closed in~\(X\).
\def\ProblemHeadingText{Problem 6.}
\item
Suppose \((X,\mathcal M,\mu)\) is a~\(\sigma\)-finite measure space. Prove the set of simple functions in \(L^2(\mu)\) is dense in \(L^2(\mu)\).
\def\ProblemHeadingText{Problem 7.}
\item
State the Baire Category Theorem. Suppose~\(X\) is a Banach space. Prove that, as a vector space, \(X\) does not have a countable basis. Show there is no norm \(\|\cdot\|\) on \(c_{00}\) such that the normed vector space \((c_{00},\|\cdot\|)\) is a Banach space.
\def\ProblemHeadingText{Problem 8.}
\item
Suppose \((X,\mathcal M,\mu)\) is a~\(\sigma\)-finite measure space. Recall that the mapping \(\lambda:L^\infty(\mu)\to L^1(\mu)^*\) defined, for \(g\in L^\infty(\mu)\), by \(g\mapsto\lambda_g:L^1(\mu)\to\mathbb C\), where, for \(f\in L^1(\mu)\),
\[
\lambda_g(f)=\int_X fg\,d\mu,
\]
 is an isometric isomorphism. Prove the following special case of this duality result. Suppose \((X,\mathcal M,\mu)\) is a finite measure space. Let \(L^1_{\mathbb R}(\mu)\) denote the real Banach space of real-valued functions in \(L^1(\mu)\). Suppose \(\lambda\in L^1_{\mathbb R}(\mu)^*\) (the dual space) and \(\lambda(f)\geq0\) for each \(f\in L^1_{\mathbb R}(\mu)\) such that \(f\geq0\) (pointwise). Prove there is an integrable function \(g:X\to\mathbb R\) such that
\[
\lambda(f)=\int_X fg\,d\mu.
\]
\end{examproblems}
\end{document}
